\documentclass[a4paper]{article}

\usepackage{graphicx}
\usepackage{amsmath,amssymb}
\usepackage{minted}
\usepackage{multirow}
\usepackage{array}
\usepackage{booktabs}
\usepackage{float}
\usepackage{algorithm}
\usepackage{algpseudocode}
\usepackage{todonotes}
\usepackage{forest}
\usepackage{xstring}
\usepackage{xcolor}
\usepackage[most]{tcolorbox}
\usepackage{xparse}
\usepackage{natbib}

\usetikzlibrary{graphs,graphdrawing}
\usegdlibrary{layered}

% for external graphviz
\input{/home/donkarlo/Dropbox/repo/nd_ai_project/src/nd_ai/knoweldge_representation/language/formal/computer/programming/tex/latex/code_clips/external_graphviz.tex}

% for tags
\input{/home/donkarlo/Dropbox/repo/nd_ai_project/src/nd_ai/knoweldge_representation/language/formal/computer/programming/tex/latex/parts/control_sequence/kind/semtag/semtag.tex}

% for links
\input{/home/donkarlo/Dropbox/repo/nd_ai_project/src/nd_ai/knoweldge_representation/language/formal/computer/programming/tex/latex/code_clips/seehere.tex}

\title{}
\date{}
\author{Mohammad Rahmani}

\begin{document}

    \maketitle

    A recurring problem in sequence modelling is to learn a model from a collection of valid training sequences and then evaluate whether a new subsequence is compatible with the process that generated those sequences. This is different from distance-based motif discovery. A probabilistic sequence model defines a probability mass function or probability density over sequences, and therefore can assign a likelihood to a test subsequence. State-space models such as Hidden Markov Models (HMMs), probabilistic suffix-tree models, and autoregressive neural models are examples of this principle \cite{murphy-2002-dynamic-bayesian-networks-representation-inference-and-learning,pimentel-clifton-clifton-tarassenko-2014-a-review-of-novelty-detection,bengio-2003-a-neural-probabilistic-language-model}.

    \section{Problem formulation}

        Suppose the training set contains sequences generated by the process of interest,

        \begin{align}
            \mathcal{D}_{\mathrm{real}}
            &=
            \left\{
            X^{(1)},X^{(2)},\dots,X^{(N)}
            \right\}.
        \end{align}

        \noindent
        A probabilistic model with parameters \(\theta_{\mathrm{real}}\) is learned from these sequences. For a test subsequence

        \begin{align}
            S
            &=
            (s_1,s_2,\dots,s_L),
        \end{align}

        \noindent
        the most direct quantity produced by a generative sequence model is

        \begin{align}
            p(S\mid\theta_{\mathrm{real}}).
        \end{align}

        \noindent
        This is the likelihood of the observed subsequence under the learned model. For discrete observations it is a probability mass. For continuous observations it is a probability density; the probability of observing one exact real-valued vector sequence is not the quantity that should be interpreted directly.

        Since raw sequence likelihood usually decreases with sequence length, a useful comparison score is the average log-likelihood,

        \begin{align}
            \ell(S)
            &=
            \frac{1}{L}
            \log p(S\mid\theta_{\mathrm{real}}).
        \end{align}

        \noindent
        Larger values indicate that the subsequence is better explained by the learned process, while smaller values indicate that it is less compatible with the training distribution.

    \section{Hidden Markov Model}

        HMMs are probabilistic state-space models for sequential data. They represent an observation sequence using an unobserved state sequence and explicitly specify an initial-state distribution, transition probabilities, and observation probabilities or densities \cite{murphy-2002-dynamic-bayesian-networks-representation-inference-and-learning}. They are therefore able to assign a likelihood to a complete test subsequence rather than only a distance or similarity score.

        Let the hidden state at time \(t\) be \(z_t\in\{1,\dots,K\}\). The HMM parameters are

        \begin{align}
            \theta
            &=
            (\boldsymbol{\pi},A,B),
        \end{align}

        \noindent
        where \(\boldsymbol{\pi}\) is the initial-state distribution, \(A\) is the transition model, and \(B\) denotes the emission model. In particular,

        \begin{align}
            \pi_i
            &=
            p(z_1=i),\\
            A_{ij}
            &=
            p(z_t=j\mid z_{t-1}=i),\\
            b_j(s_t)
            &=
            p(s_t\mid z_t=j).
        \end{align}

        \noindent
        The probability of an observed subsequence is obtained by marginalising all possible hidden-state paths,

        \begin{align}
            p(S\mid\theta)
            &=
            \sum_{z_{1:L}}
            p(z_1)
            p(s_1\mid z_1)
            \prod_{t=2}^{L}
            p(z_t\mid z_{t-1})
            p(s_t\mid z_t).
        \end{align}

        \noindent
        This quantity can be evaluated efficiently by the forward recursion. Define

        \begin{align}
            \alpha_t(j)
            &=
            p(s_{1:t},z_t=j\mid\theta).
        \end{align}

        \noindent
        The recursion is

        \begin{align}
            \alpha_1(j)
            &=
            \pi_j b_j(s_1),\\
            \alpha_t(j)
            &=
            b_j(s_t)
            \sum_{i=1}^{K}
            \alpha_{t-1}(i)A_{ij},\\
            p(S\mid\theta)
            &=
            \sum_{j=1}^{K}\alpha_L(j).
        \end{align}

        \noindent
        For continuous time-series observations, an emission can for example be Gaussian,

        \begin{align}
            p(s_t\mid z_t=j)
            &=
            \mathcal{N}
            \left(
            s_t;\boldsymbol{\mu}_j,\Sigma_j
            \right).
        \end{align}

        \noindent
        After learning the model from valid sequences, the normalized log-likelihood

        \begin{align}
            \ell_{\mathrm{HMM}}(S)
            &=
            \frac{1}{L}\log p(S\mid\theta)
        \end{align}

        \noindent
        provides a direct probabilistic compatibility score for a new subsequence.

    \section{Hidden semi-Markov model}

        A Hidden semi-Markov model (HSMM), also called a variable-duration HMM in this context, extends an HMM by explicitly modelling how long the process remains in a hidden state. This is useful when a meaningful temporal pattern consists of segments whose durations are not well described by the implicit geometric duration distribution of a standard HMM. Variable-duration semi-Markov HMMs are treated as explicit extensions of HMMs in probabilistic sequential modelling \cite{murphy-2002-dynamic-bayesian-networks-representation-inference-and-learning}.

        Let \(d\) denote the duration of a segment and let

        \begin{align}
            p(d\mid z=k)
            &=
            g_k(d)
        \end{align}

        \noindent
        be the duration distribution for hidden state \(k\). A segment ending at time \(t\), with state \(k\) and duration \(d\), explains the observations

        \begin{align}
            s_{t-d+1:t}
            &=
            (s_{t-d+1},\dots,s_t).
        \end{align}

        \noindent
        A segment-level contribution contains both a duration probability and the likelihood of the observations produced during that segment,

        \begin{align}
            p(s_{t-d+1:t},d\mid z=k)
            &=
            g_k(d)
            p(s_{t-d+1:t}\mid z=k,d).
        \end{align}

        \noindent
        The likelihood of the whole subsequence is computed by summing over possible segmentations, state sequences, and durations. Conceptually,

        \begin{align}
            p(S\mid\theta)
            &=
            \sum_{\text{valid segmentations }\mathcal{G}}
            p(S,\mathcal{G}\mid\theta).
        \end{align}

        \noindent
        Thus, as with an HMM, the HSMM produces a sequence likelihood. Its advantage is that the probability of a candidate subsequence depends not only on which states explain it but also on whether the durations of its temporal segments are plausible under the training data.

    \section{Variable-order Markov model and probabilistic suffix tree}

        A Variable-Order Markov Model (VOMM) does not impose one fixed history length on every prediction. Instead, the relevant context can be short in some parts of the sequence and longer in others. A Probabilistic Suffix Tree (PST) is a tree representation in which suffixes of different lengths define conditional probability distributions. Probabilistic suffix trees have been used to estimate conditional probability distributions from training sequences \cite{pimentel-clifton-clifton-tarassenko-2014-a-review-of-novelty-detection}.

        Assume that observations are discrete symbols from an alphabet \(\mathcal{A}\). For each position \(t\), the model selects a context

        \begin{align}
            c_t
            &=
            c(s_{1:t-1}),
        \end{align}

        \noindent
        where \(c_t\) is typically the longest suffix of the history that is represented reliably by the learned tree. The next-symbol distribution is then

        \begin{align}
            p(s_t\mid s_{1:t-1})
            &\approx
            p(s_t\mid c_t).
        \end{align}

        \noindent
        Consequently, the probability of a test subsequence can be factorized as

        \begin{align}
            p(S\mid\theta)
            &=
            \prod_{t=1}^{L}
            p(s_t\mid c_t;\theta).
        \end{align}

        \noindent
        Its average log-likelihood is

        \begin{align}
            \ell_{\mathrm{PST}}(S)
            &=
            \frac{1}{L}
            \sum_{t=1}^{L}
            \log p(s_t\mid c_t;\theta).
        \end{align}

        \noindent
        This model is especially natural when the sequence has already been discretized into symbols, states, events, words, or cluster labels. It gives an explicit probability to a discrete subsequence and automatically uses different memory lengths for different contexts.

    \section{Autoregressive probabilistic neural model and Transformer}

        Neural probabilistic sequence models estimate the distribution of the next element conditioned on previous elements. Neural language models established the use of learned distributed representations to estimate such conditional probability distributions \cite{bengio-2003-a-neural-probabilistic-language-model}. Transformer decoders provide a scalable architecture for conditioning each output on earlier positions by masked self-attention \cite{vaswani-2017-attention-is-all-you-need}.

        By the chain rule, any sequence distribution can be written as

        \begin{align}
            p(S\mid\theta)
            &=
            \prod_{t=1}^{L}
            p(s_t\mid s_{1:t-1};\theta).
        \end{align}

        \noindent
        For a discrete sequence, the Transformer can output a categorical distribution over the next symbol,

        \begin{align}
            p(s_t\mid s_{1:t-1};\theta)
            &=
            \operatorname{Categorical}
            \left(
            \boldsymbol{\pi}_\theta(s_{1:t-1})
            \right).
        \end{align}

        \noindent
        For a continuous time series, the neural network must parameterize a probability density rather than merely output a point prediction. One simple choice is a Gaussian density,

        \begin{align}
            p(s_t\mid s_{1:t-1};\theta)
            &=
            \mathcal{N}
            \left(
            s_t;
            \boldsymbol{\mu}_\theta(s_{1:t-1}),
            \Sigma_\theta(s_{1:t-1})
            \right).
        \end{align}

        \noindent
        The normalized log-likelihood of a candidate subsequence is therefore

        \begin{align}
            \ell_{\mathrm{AR}}(S)
            &=
            \frac{1}{L}
            \sum_{t=1}^{L}
            \log p(s_t\mid s_{1:t-1};\theta).
        \end{align}

        \noindent
        A high value means that the local transitions and longer-range dependencies of the candidate sequence are similar to those learned from the training sequences. A low value means that the model repeatedly assigns low probability density to elements of the candidate subsequence.

    \section{From sequence likelihood to probability of being real}

        The quantities above answer

        \begin{align}
            p(S\mid\mathrm{real}),
        \end{align}

        \noindent
        meaning ``how likely is this subsequence under the model of real training sequences?'' This is not automatically the same as

        \begin{align}
            p(\mathrm{real}\mid S),
        \end{align}

        \noindent
        meaning ``given this subsequence, what is the probability that it came from the real process?'' The distinction is important in probabilistic novelty detection, where a model of normal training data is used to decide whether test data remain compatible with normality \cite{pimentel-clifton-clifton-tarassenko-2014-a-review-of-novelty-detection}.

        If both a real-process model and a background model are available, Bayes' rule gives

        \begin{align}
            p(\mathrm{real}\mid S)
            &=
            \frac{
            p(S\mid\mathrm{real})p(\mathrm{real})
            }{
            p(S\mid\mathrm{real})p(\mathrm{real})
            +
            p(S\mid\mathrm{background})p(\mathrm{background})
            }.
        \end{align}

        \noindent
        Equivalently, a useful decision statistic is the log-likelihood ratio,

        \begin{align}
            \Lambda(S)
            &=
            \log p(S\mid\mathrm{real})
            -
            \log p(S\mid\mathrm{background}).
        \end{align}

        \noindent
        Therefore, if the desired output is literally a calibrated number such as ``this subsequence has probability \(0.87\) of belonging to the real process'', a posterior model or calibration stage is required. A single generative model by itself provides the likelihood \(p(S\mid\mathrm{real})\), not necessarily the posterior \(p(\mathrm{real}\mid S)\).

    \section{Comparison}

        \begin{table}[H]
            \centering
            \begin{tabular}{
                >{\raggedright\arraybackslash}p{0.19\textwidth}
                >{\raggedright\arraybackslash}p{0.18\textwidth}
                >{\raggedright\arraybackslash}p{0.20\textwidth}
                >{\raggedright\arraybackslash}p{0.31\textwidth}
                }
                \toprule
                Method & Typical data & Probabilistic output & Main advantage \\
                \midrule
                HMM
                & Discrete or continuous
                & \(p(S\mid\theta)\)
                & Compact latent-state model; efficient forward likelihood \\

                HSMM
                & Discrete or continuous
                & \(p(S\mid\theta)\)
                & Explicitly models the duration of temporal regimes \\

                VOMM / PST
                & Primarily discrete
                & \(p(S\mid\theta)\)
                & Variable memory length; interpretable conditional probabilities \\

                Autoregressive Transformer
                & Discrete or continuous with a density head
                & \(p(S\mid\theta)\)
                & Flexible long-range context and high-capacity modelling \\
                \bottomrule
            \end{tabular}
            \caption{Probabilistic methods for assigning likelihood to a candidate subsequence.}
        \end{table}

    \section{Practical selection}

        If the observations are continuous time-series vectors and the amount of training data is moderate, a Gaussian HMM is a strong baseline because it provides a tractable sequence likelihood and a latent representation of temporal regimes. If the duration of each regime itself is informative, an HSMM is preferable. If the observations are discrete symbols and variable context length is important, a probabilistic suffix tree is a natural choice. If a large training set is available and long-range nonlinear dependencies are important, an autoregressive Transformer with a proper probabilistic output distribution is the most flexible of these alternatives.

        For the specific goal of deciding whether a test subsequence could have been generated by the same process as the training sequences, the minimum required output is

        \begin{align}
            \boxed{
            \ell(S)
            =
            \frac{1}{L}\log p(S\mid\theta_{\mathrm{real}})
            }.
        \end{align}

        \noindent
        If an actual posterior membership probability is required, the preferred target is instead

        \begin{align}
            \boxed{
            p(\mathrm{real}\mid S)
            }.
        \end{align}

        \noindent
        The latter requires either a competing background model or a separately calibrated discriminative mapping from likelihood-based features to membership probability.

        \bibliographystyle{plainnat}
        \bibliography{/home/donkarlo/Dropbox/repo/nd_shared_research/bibliography/bibliography.bib}

\end{document}